\subsection{Large corners in the sixteen-arrow chain search}

\begin{lemma}[Two blocks with zero second cohomology]
\label{C:two-block-zero-H2-products}
Let a finite category $C$ have its objects partitioned into two
nonempty full subcategories $A,D$, with no arrows from $D$ to $A$.
If $H^2(BA;\mathbb Q)=H^2(BD;\mathbb Q)=0$, every product of two
classes in $H^2(BC;\mathbb Q)$ is zero.
\end{lemma}
\begin{proof}
The arrows from $A$ to $D$ form a profunctor.  Its two-sided
bar realization $E$ has augmentation maps to $BA$ and $BD$.
A nerve simplex of $C$ either stays in one block or has exactly
one crossing from $A$ to $D$.  Splitting at that crossing gives
the correspondence homotopy pushout
\[
 BC\simeq BA\mathop{\cup}^{h}_{E}BD.
\]
This is the same normalized two-sided bar construction as for
monoid blocks; the identities and intermediate object indices
are retained in the two block bars.  Collapsing the two vertex
spaces separately to two points gives an actual map of these
homotopy pushouts $q:BC\to\Sigma E$, where the suspension is
unreduced.  The cohomology Mayer--Vietoris sequence and the two
vanishing hypotheses show that every degree-two class on $BC$
is in the image of the connecting map from $H^1(E;\mathbb Q)$.
Naturality identifies that connecting map with $q^*$ on
$H^2(\Sigma E;\mathbb Q)$.  It is therefore surjective.
Positive cohomology products on a suspension vanish; pulling
back proves the assertion.  If $E$ is disconnected, its
components introduce only the usual additional degree-one
suspension classes and do not change this argument.
\end{proof}

\begin{lemma}[Two objects and five arrows]
\label{C:two-object-five-H2-zero}
Every connected two-object category with at most five arrows
has zero second rational cohomology.
\end{lemma}
\begin{proof}
If arrows exist in both directions, one of the two corners
has order one: otherwise two corner identities and at least
two additional corner arrows, together with two cross arrows,
would already require six arrows.  That object is a retract
of the other.  Deleting it leaves a corner of order at most
two, whose rational cohomology is concentrated in degree zero.

Otherwise the category is triangular.  With one cross arrow
its classifying space is the join of its two connected corner
classifying spaces.  Such a join has zero second rational
cohomology.  With at least two cross arrows the corner orders
sum to at most three.  Both rational corner algebras are
semisimple, and the relative normalized bar has length one.
Its rational cohomology in degree two is consequently zero.
\end{proof}

\begin{proposition}[The only large-corner inputs for sixteen-arrow chains]
\label{C:sixteen-linear-large-corners}
After the singleton-join, singleton-cofinal, fixed-arrow and
triangle-compression reductions in the fifteen-arrow chain
proof, every remaining three-object linear sixteen-arrow
profile has corner orders at most eight.  A middle corner of
order eight is excluded structurally.  At an extreme corner
of order eight it suffices to enumerate the complete catalog
of monoids with possibly nonzero second rational cohomology.
No other eight-element monoid is needed.
\end{proposition}
\begin{proof}
Write the corner orders as $(m,n,p)$ and the three forward
hom-set sizes as $(x,y,z)$, in directions $01,12,02$.
These are positive integers with sum sixteen.  The profile
calculation retained in
\path{code/outputs/c_sixteen_linear_scope_preliminary.json}
forms all $\binom{15}{5}=3003$ positive six-tuples directly,
then applies exactly the stated inequalities.  Of these,
$506$ are singleton joins, $430$ singleton-cofinal cases,
$189$ fixed-arrow cases, and $23$ triangle-compression cases.
The remaining $1855$ tuples have corner orders at most eight.
One of them is $(1,8,1;2,2,2)$, excluded by
Corollary~\ref{C:middle-double-homs-rational}.
The other middle-eight tuples are
\[
 (1,8,2;2,1,2),\quad (1,8,2;2,2,1),\quad
 (2,8,2;1,1,2)
\]
and the opposites of the first two.

At an endpoint of order two, both constant rational modules
are projective.  Passing to actual cross-act orbits gives
a functor inducing a rational cohomology-ring equivalence,
by the exact coinvariant comparison already proved for
triangular blocks.  Collapse each such endpoint.  The two
endpoints are now trivial, and every forward hom-set has
size at most two.  If either adjacent hom-set is a singleton,
singleton-cofinal deletion leaves a two-object category of
two trivial corners, hence a graph.  If both adjacent sizes
are two and the long size is one, triangle compression gives
a monoid with a two-by-two minimal band.  Its degree-two
rational products vanish by Proposition~\ref{B:rank-one-square}, without
any homological boundedness assumption.  If all three sizes
are two, Corollary~\ref{C:middle-double-homs-rational} applies.
These alternatives exclude every large middle corner after
the rational orbit comparison.

Now suppose the eight-element corner is at the source.
Unless both outgoing hom-sets are singletons, a previously
removed join, they consume at least three arrows.  The
complementary two-object full subcategory consequently has
at most $16-8-3=5$ arrows.  It is connected, since its forward
hom-set is nonempty, and Lemma~\ref{C:two-object-five-H2-zero}
applies.  If the eight-element corner also has zero second
rational cohomology, Lemma~\ref{C:two-block-zero-H2-products}
excludes the whole category.  Opposition proves the sink case.
Thus only the already complete possible-positive-$H^2$
eight-element catalog is required.  The generator retains
all its literal tables and both orientations; it makes no
claim that this smaller catalog contains every monoid of
order eight.
\end{proof}
